A base class can leave a function to its subclasses.
\codehl{FrameSerializable}, in \codehl{unito26.lob.frames},
is the base class of the objects that the package stores as tables.
It declares three functions under the decorator \codehl{abstractmethod},
with a description for a body:
the schema of the table, the conversion to a table, and the conversion from one.
And it defines the forms that follow from those three,
the list of records and the JSON text, in terms of them.
Listing \ref{listing.frameSerializable} shows the three declarations
and two of the derived forms.
A class with an abstract method is an \emph{abstract class}.
It cannot be instantiated,
and neither can a subclass until it has overridden every abstract method.

\begin{lstlisting}[caption={\protect\codehl{unito26.lob.frames.FrameSerializable}, abridged: three declarations, and two of the forms derived from them.}, label=listing.frameSerializable]
class FrameSerializable(ABC):
    # ...
    @classmethod
    @abstractmethod
    def schema(cls) -> pa.DataFrameSchema:
        # ...
    @abstractmethod
    def to_frame(self) -> pd.DataFrame:
        """This object as a validated frame."""

    @classmethod
    @abstractmethod
    def from_frame(cls, frame: pd.DataFrame) -> "FrameSerializable":
        """Rebuild from a frame, validating it first."""
    # ...
    @classmethod
    def from_records(cls, records: list[dict]) -> "FrameSerializable":
        return cls.from_frame(pd.DataFrame(records))
    # ...
    @classmethod
    def from_json(cls, text: str) -> "FrameSerializable":
        return cls.from_records(json.loads(text))
\end{lstlisting}

\codehl{MarkParams} and \codehl{HawkesParams} are subclasses, and each defines the three.
The call \codehl{MarkParams.from\_json(text)} finds \codehl{from\_json} in the base class,
and by the third row of Table \ref{tab.kindsOfMethod} the function receives
\codehl{MarkParams} as \codehl{cls}.
The references \codehl{cls.from\_records} and \codehl{cls.from\_frame} that follow
are therefore resolved along the linearisation of \codehl{MarkParams},
as in Proposition \ref{prop.lateBinding},
and the result is an instance of \codehl{MarkParams}.
This is the habit of Section \ref{sec.learningPython},
that a type carries the schema that validates it
and the pair of functions that write it to a table and read it back,
written once for every type that has it.
Listing \ref{listing.roundTrip} closes the round trip.

\begin{lstlisting}[caption={A form inherited from the abstract class returns an instance of the subclass.}, label=listing.roundTrip]
import pytest

from unito26.lob.frames import FrameSerializable
from unito26.lob.simulate import MarkParams

marks = MarkParams(
    depth_decay=0.5, mean_log_size=4.0, sigma_log_size=0.75, lot=10
)
again = MarkParams.from_json(marks.to_json(2))
assert again == marks and again is not marks
assert type(again) is MarkParams
assert "from_json" not in vars(MarkParams)

with pytest.raises(TypeError):
    FrameSerializable()
\end{lstlisting}

\begin{remark}\label{remark.twoFormsOfAbstract}
 The package leaves a function to its subclasses in a second way.
 \codehl{\_HawkesState} holds the state of a simulated process
 and the loop \codehl{events} that yields its events one \codehl{step} at a time,
 and its own \codehl{step} raises \codehl{NotImplementedError}.
 Each of its two subclasses binds \codehl{step} alone:
 \codehl{OgataThinningHawkes} is Algorithm \ref{algo.ogata},
 and \codehl{ExponentialHawkes} is Algorithm \ref{algo.dassiosZhao}.
 Nothing prevents an instance of this base class from being created,
 and the error is raised at its first event.
\end{remark}

Definition \ref{def.attributeReference} asks whether a name is found,
and never what the class of the object is.
So a function accepts as an argument any object
on which the attribute references of its body succeed, whatever the class of the object.
This is \emph{duck typing},
and inheriting from a common base class is one way of meeting it, not the only one.
\codehl{IdealisedFlow}, which places the events of a process as orders
under Assumption \ref{assumption.idealisedBook},
reads two attributes of its source of events:
\codehl{params}, for the number of types, and \codehl{events}.
Its constructor annotates the source as an \codehl{ExponentialHawkes},
and an annotation is not checked when the program runs.
The tests of the package hand it a class that yields a script of events,
so as to place one chosen event at a time,
and Listing \ref{listing.scriptedSource} does the same:
a limit buy that joins the best bid,
and a market buy that takes one share from the best ask.
The attributes that a function requires in this way can be written down
as a \codehl{typing.Protocol}, for a static type checker to verify
\citep[Chapter 13]{Ram22flu};
the package states none.

\begin{lstlisting}[caption={A scripted source of events in the place of the process.}, label=listing.scriptedSource]
from unito26.lob import config
from unito26.lob.hawkes import ExponentialHawkes
from unito26.lob.idealised import (
    IdealisedFlow, UniformDepth, idealised_book,
)
from unito26.lob.messages import BUY, SELL, GridDepth
from unito26.lob.simulate import EventType


class Scripted:
    params = config.resilient_order_flow_params()

    def __init__(self, script):
        self.script = script

    def events(self, horizon):
        return iter(self.script)


script = [(1.0, EventType.LIMIT_BUY), (2.0, EventType.MARKET_BUY)]
source = Scripted(script)
assert not isinstance(source, ExponentialHawkes)

depth = UniformDepth(5)
touches = {BUY: (1000, 2), SELL: (1010, 4)}
book = idealised_book(depth, touches, GridDepth(3))
for message in IdealisedFlow(source, depth).stream(book, 10.0, None):
    book.apply(message, record=False)
assert (book.best_bid_size, book.best_ask_size) == (3, 3)
\end{lstlisting}
